% ═══════════════════════════════════════════════════════════════════
%  Section 7 — Discussion
% ═══════════════════════════════════════════════════════════════════

\label{sec:discussion}

This section interprets the experimental results, identifies the principal methodological tensions, situates the findings relative to prior work, and identifies the limitations of the current study.

% ───────────────────────────────────────────────────────────────────
\subsection{What the Simulator Successfully Demonstrates}
\label{sec:disc-positive}

The experimental programme reveals a consistent set of outcomes that the platform can credibly demonstrate: market-level regularities that are structurally determined by the negotiation protocol, private constraints, and the feasibility gate.

On the structural side, agents reliably produce ZOPA-determined deal success (Experiment~B), stable market prices with persistent within-tick dispersion (Experiment~D), first-mover surplus advantage (Experiments~A and~D), directionally correct responses to supply and demand shocks (Experiments~E and~I), and mechanism-driven efficiency gains (Experiment~H).  These outcomes depend primarily on the existence and width of the zone of possible agreement, the alternating-offer protocol, and the matching mechanism---not on whether the agent reasons strategically.  A trivially rational agent that accepts any positive-surplus offer would produce qualitatively similar results in wide-ZOPA scenarios, though the specific price paths would differ.  The gate, which implements exactly this logic, drives approximately 40\% of all acceptances.

Experiment~H demonstrates that the matching mechanism---an institutional design choice external to the agents---has a substantial effect on market outcomes.  Surplus-maximising matching improves deal rates by 16 percentage points and allocative efficiency by 5.8 percentage points relative to random matching, while also reducing price dispersion.  These gains arise entirely from pair selection: the mechanism filters out low-ZOPA pairs before negotiation begins, ensuring that LLM agents face problems they can solve.  The surplus-maximising matcher is an oracle mechanism---it requires knowledge of both parties' private constraints---but it establishes an upper bound on the efficiency gains achievable through mechanism design in LLM-agent markets.

The asymmetric surplus distribution under surplus-maximising matching---buyer surplus rises 38\% while seller surplus is essentially unchanged---suggests that mechanism design interacts with the first-mover protocol.  Wider ZOPAs created by better matching amplify the buyer's first-mover advantage rather than distributing surplus equally.  This has implications for mechanism designers seeking both efficiency and equity in automated negotiation markets.

Experiment~I provides the clearest supply-side confirmation: a +\$20 cost shift produces near-complete price pass-through (+17.9\%), replicated across three seeds.  The welfare and surplus redistribution results are directionally consistent with supply-demand theory.  Under supply shock, the buyer/seller surplus ratio narrows from 1.69:1 (\$30.03 vs.\ \$17.77 at baseline) to 1.04:1 (\$20.79 vs.\ \$19.98), consistent with sellers passing through higher costs under a compressed ZOPA.

Taken together, these positive findings indicate that the platform can serve as a controlled testbed for evaluating how protocol and institutional structure shape market outcomes, even when individual agent reasoning is limited.

% ───────────────────────────────────────────────────────────────────
\subsection{Why Some Experiments Fail or Become Ambiguous}
\label{sec:disc-failure}

Several experiments produce null results or results too confounded to interpret cleanly.  These failures share a common source: design decisions that were engineering necessities introduced confounds that prevent clean causal attribution.

\paragraph{The gate as confound.}
The feasibility gate fires at round~2 in every fixed-scenario session, limiting concession data to two observations per session and making settlement timing invariant to deadline manipulation.  This is the primary reason Experiment~C produces a null result: deals settle before deadline proximity becomes relevant.  Extended concession trajectories are unobservable, and 100\% deal rates in wide-ZOPA scenarios may reflect gate capability rather than agent capability.  Without gate ablation, agent capability cannot be separated from gate enforcement.

\paragraph{Prompt steering confound.}
The prompt instruction to prefer \textsc{counter} over \textsc{reject} removes the primary mechanism by which agents could delay settlement.  In combination with the gate, this instruction makes deadline sensitivity structurally impossible to observe in the current design.  The instruction was introduced for a practical reason: without it, small language models frequently reject profitable offers, producing meaningless 0\% deal rates.  This illustrates a recurring tension between engineering reliability and experimental cleanliness.

\paragraph{The anchoring null.}
The accept-condition block pre-computes the utility of the opponent's standing offer and presents it directly in the prompt.  In human anchoring experiments~\cite{galinsky2001}, the anchoring effect depends on negotiators forming impressions under uncertainty about the opponent's reservation price.  Our agents receive explicit surplus information, potentially overriding any tendency to anchor on the first offer.  This is a plausible explanation for the anchoring null result ($r = -0.148$, Experiment~B), but it has not been tested: an ablation that removes the accept-condition block would be needed to confirm or reject this hypothesis.  The independent variable in Experiment~B also simultaneously manipulates ZOPA width and first-offer level, meaning anchoring cannot be fully isolated from ZOPA feasibility effects.

\paragraph{Demand-side attenuation.}
Experiment~I identifies a striking asymmetry: symmetric \$20 shifts to buyer values and seller costs produce price responses of +2.0\% and +17.9\% respectively.  One structural explanation is that seller cost functions as a hard lower bound on acceptable prices, embedded in both the gate and the agent's prompt-level constraints, producing near-mechanical cost pass-through.  Buyer value, by contrast, is a soft upper bound that buyers do not fully exploit: in Experiment~A, reactive buyers open at \$65.46 against a \$120 value (Table~\ref{tab:concession}), and the first-mover protocol means this conservative opening anchors the negotiation downward.  A value increase raises the ceiling that buyers already do not approach, producing an attenuated price response.  Alternative explanations include ZOPA-mediated mechanical effects and model-specific prompt response patterns.  The three-seed directional replication strengthens confidence in the observed patterns but does not resolve the mechanism question.

The supply shock also reveals that LLM agent effectiveness degrades as the bargaining problem becomes harder.  Table~\ref{tab:supply-demand-decomp} shows that the deal success rate among mechanically feasible pairs falls from 70.0\% to 41.1\% under supply shock, indicating that narrow ZOPAs impair LLM convergence even when agreement is structurally possible.  This parallels the finding in Experiment~B that the model achieves only a 2.7\% deal rate when the ZOPA narrows to \$20 (Table~\ref{tab:anchoring}).  Both patterns point to the same qualitative conclusion: LLM-agent market liquidity is doubly sensitive to cost shocks, through both the mechanical reduction in feasible pairs and the behavioural degradation among those that remain.

\paragraph{Prompt architecture as cognitive architecture.}
In human negotiation research, the cognitive architecture of the bargainer is relatively fixed; experimental variation comes from the task, the opponent, and the information structure.  In LLM negotiation, the prompt \emph{is} the cognitive architecture, and any change---adding a reasoning scaffold, including or excluding surplus information, steering toward or away from rejection---can alter outcomes.  LLM negotiation results are therefore inherently prompt-conditional.  The deliberative prompt's two-step reasoning scaffold (``Is the offer acceptable?  If not, what price?'') leads the model to hold a firmer position on its first counter (\$7.21 difference, Experiment~A), consistent with Abdelnabi et al.~\cite{abdelnabi2024llmdeliberation}, but whether this reflects genuine multi-step reasoning or simply a constrained output distribution cannot be determined from the current experiments.

% ───────────────────────────────────────────────────────────────────
\subsection{Micro--Macro Tension}
\label{sec:disc-micro-macro}

A structural tension runs through the entire experimental programme: designs that preserve raw agent behaviour may weaken aggregate market regularity, while stronger enforcement can stabilise outcomes but partially replaces agent decision-making.

In the micro-level experiments (A--C), the feasibility gate is enabled and fires at round~2 in every session.  This produces 100\% deal rates in wide-ZOPA conditions and allows clean comparison across prompt conditions.  The cost is that the multi-round concession dynamics the experiments were designed to study are never observed.  Between the model's generation and the enacted action, three successive mediation layers operate: the rethink loop corrects format and self-constraint violations inside the agent, the ActionJudge enforces protocol and global-bound constraints at the session level, and the gate (when enabled) overrides the agent's decision entirely.  The aggregate outcome metric---deal rate---is stable and comparable, but it is the gate's outcome, not the agent's.

In the market-level experiments (D, E, H, I), the gate continues to enforce acceptance of clearly profitable offers.  This is sufficient to generate observable market structure---price convergence, surplus asymmetry, shock response---but it means that market liquidity is partly gate-enforced rather than emergent.  The counterfactual without the gate can be inferred from Experiment~B's narrow-ZOPA condition, where the gate cannot fire and the model fails 97.3\% of the time.  A gate-disabled market experiment would likely produce near-zero deal rates with a 3B-parameter model, making market-level analysis impossible.

This reveals a fundamental constraint on LLM simulation research with small models: the researcher must choose between experimental cleanliness (gate off, agent behaviour is unmediated, but the market may not function) and practical viability (gate on, market functions, but agent behaviour is partially replaced by enforcement logic).  There is no design that fully satisfies both requirements at once for a 3B-parameter model in alternating-offer bargaining.

The tension is not unique to this platform.  Any structured economic simulation using LLM agents must decide how much enforcement to impose to make the experiment function.  The contribution of the current work is to make this trade-off explicit, quantify the gate's involvement rate (approximately 40\% of acceptances globally), and log every gate-fired settlement separately---allowing readers to assess the extent of enforcement in any given experiment.  This transparency is a prerequisite for any future comparison across model sizes, where larger models may require less enforcement and the gate involvement rate may fall.

% ───────────────────────────────────────────────────────────────────
\subsection{Comparison to Prior Work}
\label{sec:disc-prior}

We compare our findings to the prior work reviewed in Section~\ref{sec:related}.

\paragraph{Generative agent simulations.}
Park et al.~\cite{park2023generative} demonstrated that LLM agents exhibit emergent social behaviours in open-ended environments.  Our market experiments (D and~E) extend this paradigm to economic settings: LLM agents produce emergent market structure---stable price equilibrium, persistent price dispersion, and surplus asymmetry---without being explicitly programmed to produce these aggregate outcomes.  The key difference is that our environment is constrained (hard budget and cost limits, structured protocol) rather than open-ended, enabling direct comparison against theoretical predictions.

\paragraph{LLM negotiation.}
Abdelnabi et al.~\cite{abdelnabi2024llmdeliberation} find that deliberation prompting improves agreement quality in multi-party, multi-issue negotiation games.  Our Experiment~A provides a complementary finding in a simpler setting: deliberative prompting produces more conservative counter-offers (\$7.21 difference), consistent with improved strategic positioning, though we observe it through a different metric (offer price rather than Pareto efficiency).  Our setting is single-issue bilateral bargaining, which isolates price dynamics from the issue-trading effects present in their multi-issue design.

Mao et al.~\cite{mao2024personalities} show that personality-assigned prompts shift negotiation outcomes.  Our prefer-counter instruction functions similarly as a prompt-level modifier.  The platform additionally implements communication-strategy modifiers (assertive, collaborative, strategic) that are analogous to Mao et al.'s personality prompts, but these were not activated in the reported experiments; all runs used the default neutral strategy.  The methodological difference is that we disclose prompt steering as a confound rather than studying it as a controlled variable.

Lewis et al.~\cite{lewis2017dealornodeal} found that end-to-end trained negotiation agents develop non-human-like strategies, including degenerate repetitive behaviour.  Our LLM agents show a different form of departure from human norms: rather than degenerating, they fail to reproduce specific human heuristics (anchoring, deadline pressure) while maintaining otherwise reasonable bargaining behaviour.  Both findings suggest that artificial agents do not spontaneously replicate the full spectrum of human negotiation behaviour, though the underlying mechanisms differ---reinforcement learning via self-play versus prompt-conditioned inference.

\paragraph{Classical theory.}
Our positive results confirm several qualitative predictions from bilateral bargaining theory.  Market prices stabilise rapidly and exhibit persistent within-tick dispersion (consistent with Rubinstein and Wolinsky~\cite{rubinstein1985}), ZOPA width determines deal feasibility (consistent with Raiffa~\cite{raiffa1982}), and the buyer-first protocol produces buyer-favourable surplus splits (consistent with first-mover advantage in alternating-offer models~\cite{rubinstein1982}).  Experiment~I strengthens the supply-side confirmation with a three-seed directional replication: a +\$20 cost shift produces near-complete price pass-through (+17.9\%).  The demand-side confirmation is weaker: a symmetric +\$20 value shift produces only a +2.0\% price increase, directionally correct but substantially attenuated relative to the input perturbation.

Our null results depart from two well-established human findings.  The absence of anchoring bias contradicts Galinsky and Mussweiler~\cite{galinsky2001}, though the departure is plausibly explained by the prompt's provision of explicit surplus information---a condition that does not hold in human experiments.  The absence of deadline sensitivity contradicts Roth et al.~\cite{roth1988deadline}, though this result is heavily confounded by the feasibility gate and prompt steering, and should not be interpreted as evidence that LLM agents lack the capacity for deadline awareness under different experimental conditions.

% ───────────────────────────────────────────────────────────────────
\subsection{Limitations}
\label{sec:disc-limitations}

We identify eleven limitations of the current study, ordered by their impact on the interpretation of results.

\begin{enumerate}
    \item \textbf{Single model.}  All results are specific to \texttt{llama3.2:3b} at temperature~0.2.  Larger models, different architectures, or fine-tuned models may exhibit anchoring bias, deadline sensitivity, or other behaviours absent in our experiments.  No cross-model comparison is provided.

    \item \textbf{Gate dominance.}  In wide-ZOPA scenarios, the feasibility gate drives all acceptances.  Deal rates and settlement timing in Experiments~A--C reflect gate logic rather than agent decision-making.  Without gate ablation, agent capability cannot be separated from gate enforcement.  Furthermore, even in gate-disabled conditions, the rethink loop remains active and partially pre-filters the action space by correcting self-constraint violations before they reach the ActionJudge; gate-off configurations therefore do not expose fully unmediated LLM output.

    \item \textbf{Two-round data.}  The gate fires at round~2 in all fixed-scenario sessions, yielding only two data points per session for concession analysis.  Multi-round concession dynamics---the core phenomenon the concession experiment was designed to study---are not observed.

    \item \textbf{No prompt ablation.}  The accept-condition block, prefer-counter steering, and deadline-salience language are all active but none are individually ablated.  Their separate effects on anchoring, deal rates, and deadline sensitivity are unknown.

    \item \textbf{Single seed for market experiments.}  Experiments~D and~E use seed~42 only.  The specific agent pairings, ZOPA distributions, and shock realisations are fixed.  Results are suggestive but not confirmatory.

    \item \textbf{Confounded anchoring design.}  The independent variable in Experiment~B simultaneously manipulates ZOPA width and first-offer level.  The Pearson $r$ metric partially addresses this, but anchoring cannot be fully isolated from ZOPA feasibility effects.

    \item \textbf{No statistical tests.}  All analyses are descriptive (means, standard deviations, Pearson $r$).  No hypothesis tests, confidence intervals, or effect sizes are reported.  The three-seed design in Experiments~A--C provides limited replication; Experiments~D--E have no replication.

    \item \textbf{Single-price negotiation.}  All sessions negotiate a single price.  Real markets involve quality, delivery terms, and bundling.  Multi-issue negotiation may produce qualitatively different dynamics, including the logrolling and issue-trading effects studied by Lewis et al.~\cite{lewis2017dealornodeal}.

    \item \textbf{No cross-run learning.}  Agents do not accumulate experience across sessions within the completed experiments.  The memory and reputation agents are implemented but untested.  Whether LLM agents adapt strategies over repeated interactions remains an open question.

    \item \textbf{LLM non-determinism.}  Even at temperature~0.2, model outputs are not fully deterministic.  The seeded RNG controls market-level randomness (population generation, matching, shocks) but not LLM inference randomness.  This introduces uncontrolled variance into all LLM-agent conditions.

    \item \textbf{Demand-supply asymmetry mechanism unresolved.}  The demand-supply price asymmetry in Experiment~I (symmetric \$20 shifts producing +2.0\% versus +17.9\% price responses) is documented but not mechanistically explained.  Three candidate explanations---hard cost-floor constraints, conservative buyer first-mover positioning, and nonlinear ZOPA sensitivity---are consistent with the data but have not been isolated through targeted ablation.
\end{enumerate}
